\documentclass{article}
\usepackage{booktabs}
\usepackage{amsmath}
\usepackage{geometry}
\usepackage{caption}

\geometry{a4paper, margin=1in}

\title{Lab Course: Interaction and Correlations \\ Assignment \#1 -- Data Collection Templates}
\author{University of Bologna}
\date{2026}

\begin{document}

\maketitle

\section{Overview}
This standalone document contains the structured table templates required for reporting the findings of Assignment \#1. The text, headers, and captions are fully formatted in English using the \texttt{booktabs} package conventions (\texttt{\textbackslash toprule}, \texttt{\textbackslash midrule}, and \texttt{\textbackslash bottomrule}) to meet professional academic publishing standards. 

\section{Activity 1: Monte Carlo Estimate of $\pi$}

\subsection{Convergence of $\pi$ Estimate with Sample Size $N$}
This table is designed to track and compare the convergence performance of the Area Method against the direct MC Estimator as the total number of independent random samples $N$ increases. It evaluates how closely each method approaches the analytical value of $\pi$ and how the variance scales with sample size.

\begin{table}[htbp]
\centering
\caption{Convergence tracking of the $\pi$ estimate and numerical variance across different sample sizes ($N$).}
\label{tab:pi_convergence}
\begin{tabular}{llcc}
\toprule
Number of Samples ($N$) & Method & Mean ($\pi$ Estimate) & Variance \\
\midrule
1,000   & Area Method & & \\
        & MC Estimator & & \\
\addlinespace
10,000  & Area Method & & \\
        & MC Estimator & & \\
\addlinespace
50,000  & Area Method & & \\
        & MC Estimator & & \\
\bottomrule
\end{tabular}
\end{table}

\subsection{Statistical Comparison over $M$ Independent Runs}
By setting $N = 50,000$ and executing $M = 500$ independent trials with distinct random number sequences, we reconstruct the full distribution of estimators. This table synthesizes the global mean and the empirical variance of the sampled distributions, highlighting the structural precision advantages of direct function sampling over binary counting.

\begin{table}[htbp]
\centering
\caption{Statistical properties and variance comparison of the sampled distributions for $N = 50,000$ and $M = 500$.}
\label{tab:statistical_comparison}
\begin{tabular}{lcc}
\toprule
Method & Mean of the Distribution & Variance of the Distribution \\
\midrule
Area Method  & & \\
MC Estimator & & \\
\bottomrule
\end{tabular}
\end{table}

\newpage
\section{Activity 2: Markov Chain Monte Carlo (MCMC)}

\subsection{MCMC Convergence and Acceptance Rate}
This table monitors the evolution of the Markov Chain during the sampling of a bimodal Gaussian distribution under default parameters. It tracks the acceptance-to-rejection ratio across various lengths of the chain ($N$) alongside the Mean Squared Error (MSE) evaluated against the exact analytical distribution curve.

\begin{table}[htbp]
\centering
\caption{MCMC update statistics and Mean Squared Error (MSE) as a function of the total runtime samples ($N$).}
\label{tab:mcmc_convergence}
\begin{tabular}{lccc}
\toprule
Number of Samples ($N$) & Accepted Updates (\%) & Rejected Updates (\%) & MSE (vs Analytical) \\
\midrule
1,000   & & & \\
10,000  & & & \\
50,000  & & & \\
100,000 & & & \\
\bottomrule
\end{tabular}
\end{table}

\subsection{MCMC Parameter Exploration}
This template is dedicated to multi-variable parameter exploration. It documents how modifying the internal structure of the target bimodal Gaussian distribution (centers $x_0$, standard deviations $\sigma$) and clipping its physical boundaries ($x_{\min}$, $x_{\max}$) affects the Metropolis-Hastings acceptance rate during a fixed production run.

\begin{table}[htbp]
\centering
\caption{Metropolis acceptance rate dependencies under custom distribution profiles ($N = 50,000, N_{\text{therm}} = 5,000$).}
\label{tab:parameter_exploration}
\begin{tabular}{cccccl}
\toprule
$x_0$ (Center) & $\sigma$ (Std Dev) & $x_{\min}$ & $x_{\max}$ & Acceptance Rate (\%) & Notes / Configuration Typer \\
\midrule
Default & Default & Default & Default & & Baseline initialization \\
3.0     & 1.0     & -10.0   & 10.0    & & Symmetrical expanded domain \\
2.0     & 1.0     & -4.0    & 4.0     & & Truncated bimodal tail profile \\
        &         &         &         & & \\
\bottomrule
\end{tabular}
\end{table}

\end{document}
